\documentclass[11pt]{article}

\usepackage{amscd}                 %% Celui-ci provoque des erreurs!!!
\usepackage[utf8]{inputenc}
\usepackage[frenchb]{babel}
\usepackage{pictex, color}
\usepackage{amsmath, amssymb, amstext, amsthm}
\usepackage{enumerate}

\def\eqdef {\stackrel {\rm def}{=}}   % Egal par definition.
\def\eqas  {\stackrel {\rm p.s.}{=}}  % Egal a.s.
\def\toas  {\stackrel {\rm p.s.}{\to}}   % ---> conv. a.s.
\def\tod   {~\Rightarrow~}               % ---> conv. en distribution.
\def\eff {\mbox{\rm Eff}}
\def\iid {\mbox{\rm i.i.d.{}}}
\def\MSE {\mbox{\rm MSE}}
\def\RE  {\mbox{\rm RE}}
\def\rmis {{\rm is}}

\def\cF {\mathcal{F}}
\def\cH {\mathcal{H}}
\def\cS {\mathcal{S}}

\def\Var {\text{Var}}
\def\Cov {\text{Cov}}
\def\To {\Rightarrow}
\def\RR {\mathcal{R}}
\def\bA {\boldsymbol{A}}
\def\bL {\boldsymbol{L}}
\def\bI {\boldsymbol{I}}
\def\bmu {\boldsymbol{\mu}}
\def\bR {\boldsymbol{R}}
\def\bSigma {\boldsymbol{\Sigma}}
\def\bU {\boldsymbol{U}}
\def\bY {\boldsymbol{Y}}
\def\bX {\boldsymbol{X}}
\def\bx {\boldsymbol{x}}
\def\bxi {\boldsymbol{\xi}}
\def\bZ {\boldsymbol{Z}}
\def\bzero {\boldsymbol{0}}

\def\rmu{{\rm u}}

\def\dDUnif   {{\rm UniformeDiscr\`ete}}
\def\dBer     {{\rm Bernoulli}}
\def\dBinomial{{\rm Binomiale}}
\def\dGeo     {{\rm Géométrique}}
\def\dNegBinomial {{\rm BinomialeN\'egative}}
\def\dPoisson {{\rm Poisson}}
\def\dMultinomial {{\rm Multinomiale}}
%
\def\dUnif    {{\rm Uniforme}}
\def\dNormal  {{\rm Normale}}
\def\dLognormal {{\rm Lognormale}}
\def\dExpon   {{\rm Exponentielle}}
\def\dWeibull {{\rm Weibull}}
\def\dGamma   {{\rm Gamma}}
\def\dBeta    {{\rm Beta}}
\def\dStudent {\mbox{\rm Student-t}}
\def\dChiSquare {{\chi^2}}
\def\dFisher  {\mbox{\rm F}}
\def\dErlang  {{\rm Erlang}}
\def\dPareto  {{\rm Pareto}}
\def\dGumbel  {{\rm Gumbel}}
\def\dPearson {\mbox{\rm Pearson}}
\def\dPearsonV {\mbox{\rm Pearson-5}}
\def\dPearsonVI{\mbox{\rm Pearson-6}}
\def\dRayleigh {{\rm Rayleigh}}
\def\dLogistic {{\rm Logistique}}
\def\dCauchy   {{\rm Cauchy}}

\newtheorem{thm}{Theorem}
\newtheorem{prop}{Proposition}

\newcounter {exercisecount}[section]
\renewcommand{\theexercisecount}{\thesection.\arabic{exercisecount}}
\newenvironment {exercise}{}{}%\nlni\begingroup\refstepcounter{exercisecount}%
%\bf \theexercisecount\ \rm }{\xxxend}
\newenvironment {solution}{}{}%\nlni\begingroup\refstepcounter{exercisecount}%

\usepackage{verbatim}

\begin{document}

\begin{center}
\Large
Introduction à la programmation stochastique\\
Exercices - $2^e$ partie
\end{center}

\mbox{}

\hrule

\mbox{}

\mbox{}

Considérons le programme stochastique multi-étapes stylisé suivant (Heitsch et al., 2006),
cherchant à déterminer une stratégie optimale d'achat sous des coûts incertains $\bxi_t$,
sur les périodes $t = 1,\ldots,T$. Les décisions de montants à acheter sont donnés par $x_t$,
et nous souhaitons minimiser les coûts espérés de manière à atteindre un montant prescrit
$a$ à l'horizon $T$:
\begin{align*}
 \min_{\bx_t}\ & E\left[ \sum_{t=0}^T \bxi_t \bx_t \right],\\
\mbox{s.t. } & s_t-s_{t-1} = \bx_t,\ \forall t = 1,\ldots,T,\\
& s_0 = 0,\ s_T = a,\ \bx_t \geq 0,\ s_t \geq 0,
\end{align*}
où $s_t$ est une variable d'état  contenant le montant au temps $t$.

Nous considérons 3 étapes ($t = 0, 1, 2$), un arbre binaire, et $\xi_0 = 10$.
La valeur monte de 1 avec une probabilité 0.6, et descent de 1 avec une probabilité 0.4.

Doit-on ajouter des contraintes de nonanticipativité dans le problème ci-dessus?
Représentez l'arbre de scénarios.

\begin{thebibliography}{1}

\bibitem{HeitRomiStru06}
H.~Heitsch and W.~Römisch and C.~Strugarek, Stability of multistage stochastic
  programs, {\em SIAM Journal on Optimization}, 17(2):511--525 (2006).

\end{thebibliography}

\end{document}